\scanpage{454}

formulas which have a meaning only if we suppose given embeddings (or at least morphisms)

\diagram{
\[
(155)\quad
\left\{
\begin{tikzcd}[ampersand replacement=\&,column sep=large]
T_\rho \arrow[r,hook] \& G \\
T_\sigma \arrow[r,hook] \& G
\end{tikzcd}
\right.
\qquad
\begin{gathered}
\text{and a subgroup}\\
\begin{tikzcd}[ampersand replacement=\&,column sep=large]
T_0 \arrow[r,hook] \& G
\end{tikzcd}
\end{gathered}.
\]
}

When we have such embeddings, the formulas (154) have a meaning, and define \deleted{\unsure{all the}} (with the quantities (152) known) the quantities (153) in terms of one \deleted{\unsure{of them}}\inserted{of them}, say $\underline u$, or \deleted{\illegible{1}} $f$, or $\alpha$, or $\beta$ (but not in terms of $\alpha_0$ or of $\beta_0$).

In fact, we have a subgroup

\diagram{
\[
\begin{array}{r@{\;}c@{\;}l}
\begin{gathered}(156)\\[2em]\deleted{\text{\unsure{such}}}\\[3em]{}\end{gathered}
&\begin{tikzcd}[ampersand replacement=\&,row sep=3em]
{}\arrow[dd]\\{}\\{}
\end{tikzcd}
&\begin{gathered}
T_0\subset G\\
\text{such that}\quad
\left\{\begin{array}{l}
\begin{tikzcd}[ampersand replacement=\&,column sep=large]
T_0 \arrow[r,"{\substack{\delta_G|T_0\\\sim}}"] \& G_m
\end{tikzcd}\\[5pt]
\text{i.e.}\quad G\simeq T_0.SG\\
\hspace{6em}\text{product }1/2\text{ direct},
\end{array}\right.
\end{gathered}
\end{array}
\]
}

so that \deleted{the quantity}\inserted{for $U\in\Gamma G$,} there exists a \uline{unique} \deleted{$f\in\Gamma SG$ such that}\inserted{pair $(f,u)\in\Gamma(SG\times T_0)$ with}
\[
U=f^{-1}\underline u.
\]
hence unique $\alpha,\beta,\alpha_0,\beta_0\in\Gamma G$ satisfying (with the preceding $f,u$) (154) :
\[
(157)\quad
\alpha=\underline u\,\underline a^{-1},\quad
\beta=\underline u\,\underline b^{-1},\quad
\alpha_0=f^{-1}\alpha=f^{-1}\underline u\,\underline a^{-1},\quad
\beta_0=f^{-1}\deleted{\unsure{\alpha}}\inserted{\beta}=f^{-1}\underline u\,\underline b^{-1}.
\]
\mspage{453}

\scanpage{455}
\mspage{530}

We will moreover have
\[
\deleted{\unsure{\det\underline U=}}\det\underline u
=\det\underline U=\mu
\]
\[
\begin{aligned}
\det\alpha&=\underline{\det}(\underline u)\det(\underline a)^{-1}=\varepsilon'
&&\text{by (152')}\\
\det\beta&=\det(\underline u)\det(\underline b)^{-1}=1\\[3pt]
\text{and}\quad\det\alpha_0&=\det\alpha=\varepsilon'\\
\det\beta_0&=\det\beta=1
\end{aligned}
\]
we will thus indeed have
\[
\begin{aligned}
\alpha,\alpha_0&\in\Gamma G^{\pm1}\overset{\text{definition}}=
\{g\mid(\det g)^2=1\}\\
\beta,\beta_0&\in\Gamma SG
\end{aligned}
\]
--- but we have assumed, of course, that the \unsure{homomorphisms} $\delta_\rho,\delta_\sigma$ of (135) are \unsure{induced}, via the embeddings (155), by the determinant $\delta_G$.

To arrive at the relation \unsure{or} \unsure{of equation}
\[
(157)\qquad
\underbrace{[\alpha,\rho]}_{\xi}
=\varepsilon\underbrace{[\beta,\sigma]}_{\eta}\varepsilon_1(\mu-1),
\]
we must give ourselves sections $\rho,\sigma$ of $G$ (\deleted{\unsure{or}} more precisely probably, \inserted{resp.\textsuperscript{t}} of maximal $T_\rho$ and $T_\sigma$), and a strict unipotent $\varepsilon_1\in\Gamma SG$, for the relation (157) to make sense.
Let us replace $\alpha,\beta$ by \unsure{some|the} values $\underline u\,\underline a^{-1},\underline u\,\underline b^{-1}$, the relation becomes (using the commutativity of $T_\rho$, \deleted{\unsure{and}} which we had not really brought \unsure{into play} until now)
\[
(*)\qquad
[u,\rho]
=\varepsilon\cdot\underline u
\bigl(\underline b^{-1}\sigma(\underline b)^{\deleted{\unsure{-1}}}\bigr)^{\deleted{\unsure{-1}}}
\sigma(\underline u)^{-1}\varepsilon_1(\mu-1).
\]
\mspage{454}

\scanpage{456}

and will follow from the two conditions together
\[
(158)\quad
\left\{
\begin{aligned}
\text{a)}\quad&b^{-1}\sigma(b)=\varepsilon_\sigma(b)
&&\text{\uline{identically} if }b\in\Gamma N'_\sigma\\
&\text{i.e.}\quad b(\sigma)=\varepsilon_\sigma(b).\sigma
&&\text{identically}\\[4pt]
\text{b)}\quad&[u,\rho]=[u,\sigma]\varepsilon_1(\mu-1)
&&\text{identically in }u\in\Gamma T_0\\
&&&\text{for }\mu=\delta_{T_0}(u)\overset{\mathrm{dfn}}=\det(u)\\
&\text{i.e.}\quad\rho(u)=\varepsilon_1(1-\mu)\sigma(u)
&&\text{identically in }u
\end{aligned}
\right.
\]
which can also be written (applying $\sigma^{-1}$ to the two \unsure{sides|sides})
\[
\sigma^{-1}\rho(u)=\varepsilon_0(1-\mu).u
\qquad\text{where}\quad\varepsilon_0=\sigma^{-1}(\varepsilon_1)
\]
and is hardly conceivable unless \uline{$T_0$ normalizes $L_0$}, and unless $\sigma^{-1}\rho\in\Gamma T_0.L_0$ \deleted{and} is of the form
\[
(159)\quad
\left\{
\begin{aligned}
\text{a)}\quad&\widetilde{\sigma^{-1}\rho}=\widetilde{\varepsilon_0u_0}
\qquad u_0\in\Gamma T_0\\
\text{b)}\quad& T_0\text{ operates on }L_0\text{ by the \unsure{coaction|action} \unsure{induced}}\\
&\hspace{8em}\text{by the determinant}
\end{aligned}
\right.
\]
(There is all the \unsure{same} still some leeway, \unsure{in relation|by} to the draconian conditions envisaged in \uline{$\overline{\mathrm{VII}}$} !)

But here equation (157) came a bit like hair in soup --- there must be more convincing and demanding ways of bringing it in namely \deleted{by} in connection with formulas
\[
(160)\quad
\left\{
\begin{aligned}
\underline U(\rho)&=\operatorname{int}(\alpha_0).\rho
&&\bigl(=\alpha_0\rho(\alpha_0)^{-1}\rho\bigr)\\
\underline U(\deleted{\unsure{\rho}}\inserted{\sigma})&=
\varepsilon\operatorname{int}(\beta_0)\sigma
&&\bigl(=\beta_0\sigma(\beta_0)^{-1}\bigr)\\
\underline U(\varepsilon_0)&=\operatorname{int}(f^{-1})
\bigl(\varepsilon_0^{\deleted{\unsure{\mu}}}\unsure{(\mu)}\bigr)
\end{aligned}
\right.
\qquad
\begin{aligned}
\text{i.e.}\quad[\underline U,\rho]&=\alpha\rho(\alpha)^{-1}\\
[\underline U,\sigma]&=\varepsilon\beta\sigma(\beta)^{-1}
\end{aligned}
\]
\margin{watch the sign $\varepsilon$}

The third relation follows at once from $\underline U=f^{-1}\underline u$, and from the fact that by hypothesis \inserted{(159 b)} $\underline u\in\Gamma T_0$ operates on $L_0$ (\unsure{identically}) with the multiplier $\mu=\det(\underline u)$.
\mspage{455}

\scanpage{457}
\mspage{531}

The other two formulas are \unsure{made explicit} successively, by writing $\alpha=\underline u\,\underline a^{-1}$, $\beta=\underline u\,\underline b^{-1}$, $\underline U=f^{-1}\underline u$
\[
\alpha_0=f^{-1}\alpha,\quad\beta_0=f^{-1}\beta
\]
\[
(161)\qquad
\left\{
\begin{aligned}
\underline u(\rho)&=\operatorname{int}(\alpha)(\rho)\\
\underline u(\sigma)&=\varepsilon\operatorname{int}(\beta)(\sigma)
\end{aligned}
\right.
\qquad\text{where }f\text{ no longer appears}
\]
\[
\begin{aligned}
\text{i.e.}\qquad\underline u(\rho)&=\operatorname{int}(\underline u)(\rho)
&&\text{since }\operatorname{int}(\underline a^{-1})\rho=\rho\\
&&&\text{(as }T_\rho\text{ commutative)}\\
&\text{OK !}\\
\underline u(\sigma)&=\deleted{\illegible{1}}\inserted{\varepsilon}\operatorname{int}(u)
\underbrace{\bigl(\underline b^{-1}(\sigma)\bigr)}_{\substack{=\varepsilon.\sigma\\\text{cf. (158)(a)}}}
=\operatorname{int}(u)\sigma\quad\text{OK !}
\end{aligned}
\]

To prove (160) we have used only (158 a) and (159) (b). Let us now try to deduce from them the ``loop relation'' (157) --- it must be a formal calculation, which I must already have done a dozen times, in contexts which become \unsure{differentiated} as we go along ! We eliminate $f$ by writing the relations (160) in the equivalent form (161) and
\[
(161\ \text{bis})\quad
\underline u(\varepsilon_0)=\varepsilon_0(\mu)
\qquad\text{i.e.}\quad
\underbrace{[\underline u,\varepsilon_0]}_{u\varepsilon_0(u)^{-1}}
=\varepsilon_0(\mu-1)\qquad\deleted{\unsure{\operatorname{int}}}
\]
\[
\deleted{\text{i.e.}\quad u\varepsilon_0(u)^{-1}}
\]
\[
\begin{array}{l|ll}
\begin{aligned}[t]
\text{now}\qquad&\\
[\underline u,\varepsilon_0]&=
u\underbrace{\varepsilon_0(u)^{-1}}_{\substack{\unsure{\text{i.e.}}\ \sigma^{-1}\rho(u)^{\deleted{\unsure{-1}}}}}\\
\text{i.e.}\quad\sigma([u,\varepsilon_0])&=\sigma(u)\rho(u)^{-1}
\end{aligned}
&
\begin{aligned}[t]
\text{now}\quad&\deleted{\rho(u)=\rho u\rho^{-1}=}\\
&u\rho u^{-1}\rho^{-1}=\xi\\
&\rho(u^{-1})=u^{-1}\xi\\
&\rho(u)=\xi^{-1}u\\
\unsure{\text{of}=|\text{of }\overline{\mathrm m}}\quad&\sigma(u)=\deleted{\unsure{\varepsilon}}\eta^{-1}u
\end{aligned}
&
\begin{aligned}[t]
\xi&\overset{\mathrm{dfn}}{=}\alpha\rho(\alpha)^{-1}\\[4em]
\eta&\overset{\mathrm{dfn}}{=}\varepsilon\beta\sigma(\beta)^{-1}
\end{aligned}
\end{array}
\]
\[
\text{hence}\qquad
\sigma([u,\varepsilon_0])=
\deleted{\unsure{\xi^{-1}u\eta^{-1}u}}\eta^{-1}\xi
\]
and \unsure{(*)|($\alpha$)} becomes, on applying $\unsure{\sigma}$
\[
\eta^{-1}\xi=\varepsilon_1(\mu-1)
\qquad\text{i.e.}\quad\xi=\eta\varepsilon_1(\mu-1)
\qquad\text{i.e. (157)}
\]
\hfill OK !
\mspage{456}

\scanpage{458}

In summary :

\uline{Proposition} Let us suppose given, in addition to the data of the preceding section $7^\circ$ (page 526' and following \unsure{\}|)}, in particular (135)(136)), \deleted{sections $\rho$ of $\unsure{T_{\rho,\sigma}|T_\rho,\sigma}$} the following

\diagram{
\[
(162)\quad\left\{
\begin{array}{ll}
\text{a)}&\text{section }\rho\in\Gamma T_\rho,\quad\sigma\in\Gamma T_\sigma\quad\text{\uline{central} sections}\\[4pt]
\text{b)}&\text{Embeddings}\quad
\begin{tikzcd}[ampersand replacement=\&,column sep=large]
T_\rho \arrow[r,hook] \& G
\end{tikzcd},\quad
\begin{tikzcd}[ampersand replacement=\&,column sep=large]
N'_\sigma \arrow[r] \& G
\end{tikzcd}\\[4pt]
\text{c)}&\text{torus}\quad
\begin{tikzcd}[ampersand replacement=\&,column sep=large]
T_0 \arrow[r,hook] \& G
\end{tikzcd}\quad\text{such that }\delta_G:\quad
\begin{tikzcd}[ampersand replacement=\&,column sep=large]
\unsure{T_0} \arrow[r,"\sim"] \& G_m
\end{tikzcd}\\
&\text{i.e.}\quad G\simeq T_0.SG\quad\text{(semidirect)}\\[4pt]
\text{d)}&\varepsilon_0\in\Gamma G\quad\text{\deleted{strict} regular unipotent, defining}\\
&\text{thus}\quad
\begin{tikzcd}[ampersand replacement=\&,column sep=large]
\unsure{\lambda} \arrow[r,mapsto] \& \varepsilon_0(\lambda)
\end{tikzcd}:\quad
\begin{tikzcd}[ampersand replacement=\&,column sep=large]
G_a \arrow[r,hook] \& SG
\end{tikzcd}
\end{array}\right.
\]
}

and let us suppose \deleted{(in addition to the \uline{commutativity} of $T_\rho,T_\sigma$}
\[
\begin{aligned}
\inserted{\text{a)}\quad\delta_G|T_\rho=\delta_\rho,\quad\delta_G|N'_\sigma=\delta_\sigma}\\
\deleted{\text{a)}\quad T_0\text{ \unsure{normalizes} }\unsure{L_0}}
\end{aligned}
\]
\[
(163)\quad
\begin{aligned}
\text{b)}\quad&\operatorname{int}(\underline b)(\sigma)=\varepsilon_\sigma(\underline b)\sigma
\quad\text{in }G\text{ identically}\\
&\hspace{12em}\text{for }\underline b\in\Gamma\unsure{N'_\sigma|N^2_\sigma}\\[4pt]
\text{c)}\quad&T_0\text{ normalizes }L_0,\text{ and we have}\\
&u(\varepsilon_0(\lambda))=
\varepsilon_0\bigl(\deleted{\unsure{\delta_G(u)}}\mu\lambda\bigr)
\quad\text{where }\mu=\delta_G(u)\\[4pt]
\text{d)}\quad&\widetilde{\unsure{\sigma|\rho}}^{-1}\widetilde\rho
\in\deleted{\unsure{\widetilde T_0}}\Gamma\widetilde\varepsilon_0\widetilde L_0
\quad\text{i.e. }\exists\text{locally}^{\text{t}}\text{ section }\lambda_0\in\Gamma G_m\\
&\text{with}\quad\lambda_0\widetilde\sigma^{-1}\widetilde\rho=\varepsilon_0u_0,
\quad\text{with }u_0\in\Gamma T_0
\end{aligned}
\]
\margin{\unsure{links} $L_0,T_0$}
\margin{$\begin{array}{l|}\text{\unsure{links}}\quad L_0,T_0\\\rho,\sigma\end{array}$}

Under these conditions, we have for every system (152) satisfying the relations (152') (i.e. in fact for every \deleted{\unsure{syst}} section of $\mathcal M_{\rho,\sigma}$ !) a unique \unsure{system} (153), satisfying the relations (154) --- namely $\underline U$ can be written uniquely in the form
\[
(164)\quad\underline U=f^{-1}\underline u,\qquad f\in\Gamma SG,\quad\underline u\in\Gamma T_0
\]
and we put
\[
\begin{aligned}
(165)\quad\alpha&=\underline u\,\underline a^{-1},\quad
\beta=\underline u\,\underline b^{-1},\quad
\alpha_0=f^{-1}\unsure{\alpha|\alpha_\rho}=
\deleted{\unsure{\beta}}\inserted{f}^{-1}\underline u\,\underline a^{-1},\\
\beta_0&=f^{-1}\beta=f^{-1}\underline u\,\underline b^{-1}.
\end{aligned}
\]
\mspage{457}